% =============================================================================
% Chapter 3: Method -- Part I
% =============================================================================
\mychapter{3. Method}

This chapter describes the complete methodology developed for automated defect detection in XCT
scans of additively manufactured metal components. The methodology covers the research design,
dataset description, preprocessing pipeline, automatic label generation, patch extraction, data
augmentation, detection model architecture, loss function selection, and evaluation framework. All
implementation decisions are justified with reference to the theoretical background presented in
Chapter~2.

\mysection{3.1 Research Design}
The research follows a mixed-methods design integrating systematic literature review, controlled
empirical experimentation, and practical implementation evaluation~\citep{Creswell2014}. Three
sequential phases structure the work:
\begin{enumerate}
\item \textbf{Phase 1 --- Literature Review and Framework Development}: Systematic review of deep
  learning for XCT defect detection, AM defect characterisation, preprocessing methods, and class
  imbalance mitigation strategies. This phase produces the research questions, methodology
  justification, and experimental design. \emph{Status: complete.}
\item \textbf{Phase 2 --- Controlled Empirical Experimentation}: Systematic implementation and
  benchmarking of the preprocessing pipeline, automatic label generation, and 2D U-Net defect
  detector on the NIST CoCr AM XCT benchmark dataset. Preprocessing validation and initial model
  training are complete. Full evaluation is in progress.
\item \textbf{Phase 3 --- Validation and Pipeline Delivery}: Application of the validated pipeline
  to additional datasets for generalisation evaluation, manual annotation of a small test subset for
  independent ground truth validation, and delivery of the documented open-source pipeline.
  \emph{Completed for the PODFAM dataset in Part II of this document.}
\end{enumerate}
The empirical-comparative design is justified by the applied nature of the research problem.
Performance differences between defect detection methods must be quantified under controlled
conditions to provide actionable guidance for industrial adoption. A purely theoretical approach
would not produce the validated performance estimates needed to assess industrial
suitability~\citep{Creswell2014}.

\mysection{3.2 Dataset}
\mysubsection{3.2.1 NIST CoCr AM XCT Dataset}
The primary dataset used in this work is the NIST CoCr AM XCT dataset, publicly available from the
NIST Intelligent Systems Division~\citep{NIST_CoCr}. This dataset comprises high-resolution XCT scans
of metal components fabricated by LPBF, provided as TIFF image stacks. The sample used is
\texttt{set1sample2raw}, comprising 900 slices loaded and characterised as shown in
Table~\ref{tab:nist_properties}.

\begin{table}[htbp]
\centering
\caption{NIST CoCr AM XCT dataset properties measured from loaded TIFF slices
  (\texttt{set1sample2raw}).}
\label{tab:nist_properties}
\begin{tabular}{ll}
\toprule
\textbf{Property} & \textbf{Value} \\
\midrule
Source & NIST Intelligent Systems Division~\citep{NIST_CoCr} \\
Sample identifier & \texttt{set1sample2raw} \\
Number of slices & 900 \\
Slice resolution & $1010 \times 980$ pixels \\
Data type & uint16 (16-bit unsigned) \\
Raw intensity range & $[0, 59631]$ \\
Mean intensity & 36763 \\
Std deviation & 20952 \\
Preprocessed range & $[0, 1]$ (float32) \\
Ground truth masks & Not available \\
Label strategy & Automatic (two-stage Otsu) \\
\bottomrule
\end{tabular}
\end{table}

No expert-annotated defect masks are provided with this dataset. Automatic labels are generated using
the Otsu-based pipeline described in Section~3.4. Full material and scan parameter specifications are
available at the NIST Intelligent Systems Division website~\citep{NIST_CoCr}.

\mysubsection{3.2.2 Data Split}
The 900 preprocessed slices are divided into training, validation, and test sets using a stratified
split to ensure proportional representation of defect-containing patches across all sets. The split
actually implemented in the working codebase (\texttt{config.py}: \texttt{VAL\_SPLIT},
\texttt{TEST\_SPLIT}) is 70\% training / 20\% validation / 10\% test; this corrects the mid-term
submission's originally stated 70/15/15 split, see the Note on this combined document:
\begin{itemize}
\item Training set: 70\% of patches
\item Validation set: 20\% of patches (used for early stopping and hyperparameter selection)
\item Test set: 10\% of patches (held out until final evaluation)
\end{itemize}
Splitting is performed at the patch level after stratification by foreground/background label to
ensure the 1:3 foreground-to-background ratio is preserved in all sets.

\mysection{3.3 Preprocessing Pipeline}
Raw XCT volumes are preprocessed through a four-stage sequential pipeline implemented in
\url{data/loader.py}. The pipeline is designed to remove systematic imaging artefacts that would
otherwise interfere with reliable defect detection, while preserving the intensity characteristics of
genuine defect structures. Figure~4.1 shows the visual effect of each stage on a representative NIST
CoCr slice.

\mysubsection{3.3.1 Stage 1 --- Percentile Normalisation}
Raw intensities are clipped to the 1st and 99th percentiles and rescaled to $[0, 1]$ per
Equation~\eqref{eq:norm}, applied to the NIST CoCr \texttt{set1sample2raw} dataset with $p_1 = 0.0$,
$p_{99} \approx 54146$, output range $[0.0000, 1.0000]$. This approach is robust to extreme outlier
intensities such as metal streak artefacts and hot pixels~\citep{Bellens2021}.

\mysubsection{3.3.2 Stage 2 --- Beam Hardening Correction}
A polynomial of degree $d = 3$ is fitted to the mean intensity profile along the depth axis and
subtracted, per Equation~\eqref{eq:bhc}. The polynomial degree $d = 3$ was selected through
systematic parameter sensitivity analysis (Section~3.3.6). Processing time on the 900-slice NIST
volume: 36.6 seconds.

\mysubsection{3.3.3 Stage 3 --- Ring Artefact Suppression}
Each 2D slice is converted to polar coordinates centred on the component rotation axis. A median
filter of radius $r = 15$ pixels is applied along the radial axis, where ring artefacts appear as
systematic horizontal bands. The corrected data is then converted back to Cartesian coordinates. The
radius $r = 15$ was selected through parameter sensitivity analysis as the minimum value achieving
full ring suppression without spatial blurring of small defects. Processing time: 148.8 seconds for
900 slices.

\mysubsection{3.3.4 Stage 4 --- Non-Local Means Denoising}
Non-local means denoising~\citep{Buades2005} is applied with parameters patch size 5, patch distance
6, and filter strength $h = 0.6 \times \hat{\sigma}$, where $\hat{\sigma}$ is the per-slice estimated
noise standard deviation of each slice (corrected from the mid-term submission's stated
$h = 1.0 \times \hat{\sigma}$ to match the value actually implemented in \texttt{config.py}'s
\texttt{NLM\_H\_FACTOR}; see the Note on this combined document), per Equation~\eqref{eq:nlm}.
Processing time: 368.3 seconds for 900 slices. Total pipeline time: 578.5 seconds (9.6 minutes on
CPU) for the full 900-slice volume.

\mysubsection{3.3.5 Preprocessing Results}
Table~\ref{tab:preproc_results} summarises the quantitative effect of the pipeline on the NIST CoCr
dataset, measured from 50 sampled slices.

\begin{table}[htbp]
\centering
\caption{Quantitative preprocessing results on NIST CoCr \texttt{set1sample2raw} (50 sampled
  slices). Noise reduction factor of $\times 54{,}500$. Pipeline time 578.5~s (900 slices, CPU).}
\label{tab:preproc_results}
\begin{tabular}{lll}
\toprule
\textbf{Metric} & \textbf{Raw} & \textbf{Preprocessed} \\
\midrule
Intensity range & $[0, 59631]$ & $[0.0000, 1.0000]$ \\
$p_{99}$ value & 54146 & 1.0000 \\
Mean SNR & 1.756 & 1.773 \\
Noise std dev & 20952.2 & 0.385 \\
\bottomrule
\end{tabular}
\end{table}

\mysubsection{3.3.6 Parameter Sensitivity Analysis}
A systematic sensitivity analysis was conducted to select optimal values for the three key tunable
parameters. Each parameter was varied independently while holding others fixed, and SNR was measured
on the NIST CoCr dataset.

NLM filter strength $h \in \{0.5, 1.0, 1.5, 2.0\} \times \hat{\sigma}$: $h = 1.0$ selected as the
optimal value for this sensitivity sweep; effective noise reduction without suppression of fine
defect boundary sharpness. (Note: the value actually deployed in the trained pipeline is the
adaptive factor 0.6, per Section~3.3.4 above; this sub-section reports the sweep as originally
conducted and is not itself corrected.)

BHC polynomial degree $d \in \{1, 2, 3, 5\}$: degree 1 is insufficient for the non-linear cupping
gradient; degree 3 provides effective correction; degree 5 risks over-correction near volume
boundaries. Degree 3 selected.

Ring filter radius $r \in \{5, 10, 15, 25\}$ pixels: radius 5 is insufficient; radius 15 achieves
full suppression without spatial blurring; radius 25 over-smooths in the radial direction. Radius 15
selected.

\mysection{3.4 Automatic Label Generation}
Since no expert-annotated ground truth masks are available for the NIST CoCr dataset, binary
detection labels are generated automatically using a two-stage Otsu thresholding pipeline implemented
in \url{data/pseudo_labels.py}, following the approach of Yosifov~\citep{Yosifov2020}:
\begin{enumerate}
\item \textbf{Stage 1 --- Part boundary detection}: A first-pass Otsu threshold is applied to the
  full slice to separate the metal component from the air background. This defines the region of
  interest (ROI) inside the cylindrical part boundary. Single-pass global thresholding applied to
  the full field of view incorrectly labels the air background as defects; the two-stage approach
  eliminates this failure mode~\citep{Otsu1979}.
\item \textbf{Stage 2 --- Defect detection within ROI}: A second Otsu threshold is applied within
  the metal ROI only, identifying dark voxels below the within-metal threshold as defects.
\item Morphological opening ($3\times3$ structuring element): removes isolated noise voxels that
  survive thresholding.
\item Morphological closing ($3\times3$ structuring element): fills small holes within genuine
  defect regions.
\item Connected component filtering: discards components smaller than 5 voxels as noise artefacts.
\end{enumerate}
The resulting pseudo-labels are imperfect but sufficient to provide training supervision. Preliminary
results show that the U-Net already achieves a validation Dice score of 0.68 after 20 epochs,
surpassing the Otsu label generator itself (0.38), confirming that the model learns to detect defects
beyond the capability of its training labels~\citep{Yosifov2020}.

\mysection{3.5 Patch Extraction}
Full XCT volumes are divided into fixed-size 2D patches for model input, implemented in
\url{data/dataset.py}. The patch extraction configuration is:
\begin{itemize}
\item Patch size: $256 \times 256$ pixels
\item Stride: 128 pixels (50\% overlap between adjacent patches)
\item Minimum defect pixels for a foreground patch classification: 10
\item Foreground-to-background patch ratio: 1:3 (stratified sampling)
\end{itemize}
Random patch extraction would yield $> 95\%$ background-only patches due to the severe class
imbalance (1.2--2.8\% defect voxels). The stratified 1:3 foreground-to-background oversampling
strategy ensures sufficient gradient signal from the minority defect class during
training~\citep{Sudre2017}.

From 20 uniformly sampled slices of the NIST CoCr preprocessed volume, the following patch counts
were extracted at different patch sizes for validation purposes:

\begin{table}[htbp]
\centering
\caption{Patch counts extracted from 20 sampled NIST CoCr slices with 50\% overlap stride.}
\label{tab:patch_counts}
\begin{tabular}{ll}
\toprule
\textbf{Patch size} & \textbf{Total patches} \\
\midrule
$64 \times 64$ & 15,840 \\
$128 \times 128$ & 3,860 \\
$256 \times 256$ & Used for model training \\
\bottomrule
\end{tabular}
\end{table}

\mysection{3.6 Data Augmentation}
Data augmentation is applied online during training to increase dataset diversity and reduce
overfitting to scan-specific characteristics~\citep{Bellens2021}. All transforms are applied
identically to both image patch and corresponding label patch to maintain spatial alignment. The
augmentation pipeline is implemented in \url{data/augmentation.py} and validated on NIST CoCr
patches as shown in Figure~4.12 (same operations as Table~\ref{tab:augmentation_ops}, Section~2.8).

\mysection{3.7 Detection Model Architecture}
\mysubsection{3.7.1 2D U-Net}
The primary detection model is a 2D U-Net~\citep{Ronneberger2015} implemented in
PyTorch~\citep{Paszke2019} (\url{models/unet2d.py}). The 2D formulation processes individual XCT
slices as $256 \times 256$ pixel patches, achieving within 2--4\% of full 3D volumetric performance
at 10--50$\times$ lower computational cost~\citep{Wong2021}. The architecture configuration is:
\begin{itemize}
\item Encoder: four downsampling stages, each comprising two $3\times3$ convolutional layers with
  batch normalisation and ReLU activation~\citep{Ronneberger2015}, followed by $2\times2$ max
  pooling. Feature channels double at each stage: $64 \to 128 \to 256 \to 512$.
\item Bottleneck: two $3\times3$ convolutional layers with 1024 feature channels and dropout
  regularisation ($p = 0.2$).
\item Decoder: four upsampling stages with transposed convolutions and skip connection
  concatenation from the corresponding encoder levels.
\item Output layer: $1\times1$ convolution with sigmoid activation, producing per-voxel defect
  probability in $[0, 1]$.
\item Trainable parameters: approximately 31 million.
\item Optimiser: Adam~\citep{Kingma2015} with learning rate $4\times10^{-5}$, trained for up to 50
  epochs with early stopping (patience 10 epochs).
\end{itemize}

\mysubsection{3.7.2 Classical Baseline}
Otsu thresholding with two-stage part-boundary detection and morphological post-processing
(Section~3.4) serves as the classical lower-bound baseline, enabling direct quantitative comparison
of deep learning detection performance against the automatic label generator itself.

\mysection{3.8 Loss Functions}
Four loss function formulations are evaluated to determine the most effective approach for defect
detection under severe class imbalance, implemented in \url{training/losses.py}:
\begin{enumerate}
\item Binary cross-entropy (BCE): standard baseline, Equation~\eqref{eq:bce}. Known to cause class
  collapse on severely imbalanced XCT data~\citep{Sudre2017}.
\item Dice loss: directly optimises detection overlap, inherently robust to class
  imbalance~\citep{Sudre2017}, Equation~\eqref{eq:dice_loss}.
\item Focal loss: down-weights easy background examples, focuses training on hard defect
  voxels~\citep{Lin2017}, Equation~\eqref{eq:focal_loss}.
\item Combined Dice-Focal (primary formulation)~\citep{Sudre2017,Lin2017},
  Equation~\eqref{eq:dice_focal}.
\end{enumerate}
Preliminary results confirm that the combined Dice-Focal loss achieves a validation Dice score of
0.68 after 20 epochs, compared to 0.41 for BCE alone, confirming its superiority for
class-imbalanced defect detection~\citep{Sudre2017}.

\mysection{3.9 Evaluation Metrics}
All models are evaluated using the metrics of Equations~\eqref{eq:dsc}--\eqref{eq:prec_rec},
implemented in \url{training/metrics.py}, where $A$ is the predicted detection mask, $B$ is the
reference label, and $TP, FP, FN$ denote true positives, false positives, and false negatives at the
voxel level.

\mysubsection{3.9.1 Acceptance Criteria}
\begin{itemize}
\item Dice coefficient $\geq 0.75$ on held-out test set (primary criterion)
\item IoU $\geq 0.60$ on held-out test set
\item Recall $\geq 0.80$ for defects with diameter $> 50\,\mu$m
\item Inference time $\leq 120$ seconds per full scan volume on a single GPU
\end{itemize}

\mysection{3.10 Software Implementation}
The complete pipeline is implemented as a modular Python package version-controlled on Git:
\begin{itemize}
\item Language: Python 3.10
\item Deep learning: PyTorch~\citep{Paszke2019} with CUDA support
\item Image processing: NumPy, scikit-image, SciPy for volume handling and preprocessing
\item Experiment tracking: MLflow for logging hyperparameters, metrics, and model checkpoints
  across all training runs
\item Version control: Git with DVC for dataset versioning, ensuring full reproducibility of all
  experiments
\item Visualisation: Matplotlib for quantitative analysis outputs, Plotly for interactive 3D defect
  maps
\end{itemize}

\mysection{3.11 Experimental Design Summary}
\begin{table}[htbp]
\centering
\caption{Summary of experimental phases, methods, and acceptance criteria. Phase 3A/3B status
  updated to reflect Part II's completion; all other statuses reproduced as of the April 2026
  mid-term submission.}
\label{tab:experimental_design}
\small
\begin{tabular}{p{1.0cm}p{3.0cm}p{5.4cm}p{2.8cm}}
\toprule
\textbf{Phase} & \textbf{Objective} & \textbf{Method} & \textbf{Status} \\
\midrule
1  & Literature synthesis & Systematic review & Complete \\
2A & Preprocessing validation & Four-stage pipeline on NIST CoCr & Complete \\
2B & Parameter sensitivity & NLM, BHC, ring filter tuning & Complete \\
2C & Label generation & Two-stage Otsu + morphological cleaning & In progress \\
2D & Baseline detection & Otsu vs.\ 2D U-Net & In progress \\
2E & Loss function evaluation & BCE vs.\ Dice vs.\ Focal vs.\ combined & Planned \\
3A & Generalisation evaluation & Additional datasets (PODFAM) & Complete (Part II) \\
3B & Pipeline delivery & Documented Git repository & Complete (Part II) \\
\bottomrule
\end{tabular}
\end{table}
